\scanpage{158}
considérer $E$ comme le module enveloppant de $\mathop{\ast}A_i$, en le munissant de l'augmentation
\[
 \varepsilon:E\longrightarrow k
\]
définie par
\[
 \varepsilon(\underbrace{t_1t_2\cdots t_n}_{t})=2
\]
On notera que $\varepsilon$ n'\uline{est pas} induit par
\[
 \deleted{\unsure{\varphi}}\varepsilon_1\otimes\cdots\otimes\varepsilon_n:\bigotimes E_i\mathbin{\unsure{\text{---}|\longrightarrow}}k
\]
qui transforme $t_1t_2\cdots t_n$ non pas en \unsure{$2$}, mais en $2^n$.

Le groupe $(\mathbb Z/2\mathbb Z)^n$ opère sur $P=\bigotimes E_i$ via les $\sigma_i$, et les différents caractères définissent des sous-modules propres, dont $\bigotimes E_i$ est la somme directe (si $2$ inv !). Un car. est donné par $\alpha=(\alpha_1,\ldots,\alpha_n)\in(\mathbb Z/2\mathbb Z)^n$, et \deleted{\unsure{son}}
\[
 P_\alpha=\{x\in P\mid \sigma_1^{\deleted{\unsure{\alpha_1}}\beta_1}\otimes\cdots\otimes\sigma_n^{\deleted{\unsure{\alpha_n}}\beta_n}(x)=(-1)^{\alpha_1\beta_1+\alpha_2\beta_2\cdots+\alpha_n\beta_n}x\}
\]
\margin{$\forall$ syst. d'entiers $\beta_i$ mod \uline{$2$}}
Ceci dit, $E$ est le sous espace propre associé au caractère somme $\alpha=(1,\ldots,1)$ i.e.
\[
 E=P_{1,\ldots,1}=\{x\in P\mid(\sigma_1^{\beta_1}\deleted{\unsure{\otimes\sigma}}\otimes\cdots\otimes\sigma_n^{\beta_n})x_n)=(-1)^{\sum\beta_i}x\}
\]
% archive sheet number: 157

